\chapter{Grundlagen}
\label{chap:grundlagen}

Bevor ein \gls{gls-KI}-basiertes Telefonsystem entwickelt wird, lohnt sich der Blick auf Container-Umgebungen: Eine \gls{gls-VM}, die einen eigenen Gast-Kernel verwendet, nutzen im Gegensatz dazu Container den Host-Kernel direkt und sparen so Hypervisor-Overhead \parencite{Antonova2025}. Nach geltenden Richtlinien ist für die Virtualisierung von Systemkomponenten nicht zwingend eine vollständige \gls{gls-VM} erforderlich. Stattdessen können auch Container oder andere Isolationsmechanismen als leichtgewichtige Alternativen eingesetzt werden, solange sie ausreichende Isolation bieten und effiziente Snapshots unterstützen \parencite{ETSI-REL}. In Cloud-nativen Architekturen bilden Container und zusätzliche Orchestrierung den Standard, weil Bereitstellung, Skalierung und Resilienz damit automatisiert werden \parencite{ReliabilityVirtualizedNetworks}.

\section[Container, Isolation und
Deployment]{Container, Isolation und
Bereitstellung}

Dieser Ansatz macht Container deutlich leichter und schneller als \gls{gls-VM}. Anwendungen starten innerhalb von Millisekunden und
benötigen weniger Speicher. Gerade in heterogenen Umgebungen oder bei begrenzter Hardware lassen sich so viele Dienste effizient parallel betreiben. Dies wird auch in Arbeiten zu containerbasierten Telefonsystemen bestätigt, die zeigen, dass Container auf ressourcenschwacher Hardware meist stabiler und performanter laufen als klassische \gls{gls-VM} \parencite{Antonova2025}.
Für die Entwicklung komplexerer Systeme wie \gls{gls-VoIP}-Umgebungen oder \gls{gls-KI}-basierte Dienste sind zudem die Eigenschaften von Containerplattformen wie Docker relevant. Docker-Images bestehen aus mehreren Schichten, wobei jede Schicht eine Änderung am Dateisystem repräsentiert. Diese Schichten werden inkrementell zu einem vollständigen Image zusammengesetzt, was das Verteilen und Aktualisieren vereinfacht. Durch dieses Schichtmodell können identische Basisschichten von mehreren Images gemeinsam genutzt werden, was Speicher spart und das Caching sowie die Wiederverwendbarkeit fördert. Sicherheitsupdates oder Änderungen an einzelnen Komponenten lassen sich gezielt und effizient ausrollen, ohne das gesamte Image neu erstellen zu müssen. In der Praxis führt dieser modulare Aufbau dazu, dass Systeme flexibel erweitert, skaliert und besser gewartet werden können \parencite{docker-docs}.

Gleichzeitig muss man die Grenzen dieses Ansatzes kennen: Da Container sich den Kernel teilen, hängt ihre Sicherheit stark von dessen Stabilität ab. Fehler oder Schwachstellen im Kernel können sich auf alle laufenden Container auswirken. Diese Einschränkung spielt besonders dann eine Rolle, wenn sensible Daten wie Sprachdaten verarbeitet werden. Um eine vollständig sichere und stabile, bidirektionale Audioübertragung zu gewährleisten, müssen die Dienste in Containern einwandfrei miteinander kommunizieren und die Container fehlerfrei miteinander harmonieren \parencite{Antonova2025,ETSI-REL,ReliabilityVirtualizedNetworks}.

\vspace{1em}

\section[Klassische Telefonieprotokolle]{Klassische Telefonieprotokolle}

Um zu verstehen, wie sich \gls{gls-KI}-Komponenten in ein Telefonsystem integrieren lassen, ist ein Blick auf die technischen Grundlagen der Internet-Telefonie notwendig. \gls{gls-VoIP} stellt dabei die Standardtechnologie zur Übertragung von Sprachdaten über IP-Netzwerke dar. Der zugrunde liegende Protokollstapel basiert auf dem Zusammenspiel mehrerer standardisierter Protokolle \parencite{rfc3261,ITU-T-G114}.

Für den Rufaufbau und die Signalisierung wird hauptsächlich das \gls{gls-SIP}-Protokoll eingesetzt. Es handelt sich dabei um ein Standardprotokoll zur Steuerung von Sprach- und Videokommunikation \parencite{rfc3261}. \gls{gls-SIP} regelt unter anderem die Registrierung von Endgeräten sowie den Aufbau und Abbau von Sprachverbindungen über einen\newline INVITE/ACCEPT/BYE-Handshake sowie Registrierungsmechanismen zur Anmeldung von Clients an Proxy-Servern. Die eigentlichen Audiodaten werden anschlie\ss end über \gls{gls-RTP} übertragen, das für die Echtzeitübertragung von Audio- und Videodaten ausgelegt ist. Ergänzend liefert \gls{gls-RTCP} als Begleitprotokoll zu \gls{gls-RTP} Informationen über die Verbindungsqualität und stellt Feedback zur laufenden Sitzung bereit \parencite{rfc3550}.

Open-Source-Systeme wie Asterisk \footnote{Offizielle Webseite: \url{https://www.asterisk.org/}} oder Kamailio \footnote{Offizielle Webseite: \url{https://www.kamailio.org/w/}} setzen konsequent auf diese Standards. Diese Protokollstandards finden sich beispielsweise in einer \gls{gls-PBX} wieder. Ein \gls{gls-PBX} ist eine Telefonanlage für interne und externe Sprachkommunikation, die Funktionen wie Anrufweiterleitung und Konferenzschaltungen bereitstellt. Als bewährte Open-Source-Lösungen bieten solche Systeme vielfältige Möglichkeiten zur Anrufverarbeitung, die von einfachen Durchwahlszenarien über \gls{gls-IVR} mit automatisierten Sprachmenüs bis hin zu komplexeren Konferenzfunktionen reichen \parencite{Antonova2025}. In modernen Architekturen werden \gls{gls-PBX}-Systeme häufig in containerbasierten Umgebungen und häufig in Verbindung mit Cloud-Diensten. betrieben. Der \gls{gls-PBX}-Kern, der \gls{gls-SIP}-Proxy, eine Datenbank und gegebenenfalls ein \gls{gls-RTP}-Proxy laufen dabei als separate Dienste, wodurch sich die Skalierbarkeit und Fehlertoleranz erhöht \parencite{Antonova2025,ReliabilityVirtualizedNetworks}.

Für containerbasierte Umgebungen ist das \gls{gls-NAT}-Traversierung besonders relevant, da die Echtzeitkommunikation auf bidirektionale und dynamische Verbindungen angewiesen ist. \gls{gls-SIP} wurde ursprünglich für Netzwerke mit öffentlich erreichbaren Endpunkten konzipiert und berücksichtigt keine durch \gls{gls-NAT} verursachte Adressumsetzung \parencite{rfc3261}. Insbesondere bei der Aushandlung von Medientransporten werden IP-Adressen und Ports im Protokollinhalt übertragen, die aus der Sicht des \gls{gls-NAT} häufig nicht erreichbar sind, was den direkten Austausch von Audio- und Videodaten verhindert.

Zur Überwindung dieser Einschränkungen kommen zusätzliche \gls{gls-NAT}-Traversierungsmechanismen zum Einsatz. Das \gls{gls-STUN}-Protokoll ermöglicht Endpunkten die Ermittlung ihrer von au\ss en sichtbaren Transportadressen \parencite{rfc5389}. Das \gls{gls-TURN}-Verfahren stellt als Fallback eine relaybasierte Weiterleitung von Mediaströmen bereit, wenn direkte Verbindungen aufgrund restriktiver \gls{gls-NAT}-Typen nicht möglich sind \parencite{rfc5766}. Aufbauend darauf koordiniert \gls{gls-ICE} diese Verfahren, indem verschiedene potenzielle Verbindungspfade systematisch getestet werden, um eine funktionierende Ende-zu-Ende-Verbindung auszuhandeln \parencite{rfc5245}.

In containerbasierten Umgebungen wie Kubernetes\footnote{Offizielle Webseite: \url{https://kubernetes.io/}} oder Docker\footnote{Offizielle Webseite: \url{https://www.docker.com/}} wird diese Problematik durch zusätzliche virtuelle Netzwerkebenen weiter verschärft, da Container in der Regel über mehrere Adressübersetzungen hinweg kommunizieren müssen, was NAT-Traversierung zu einer zentralen technischen Herausforderung für die Echtzeitkommunikation in Telefonsystemen macht\parencite{Antonova2025,ReliabilityVirtualizedNetworks}.

\vspace{1em}

\section{Audio-Formate und Transkodierungs-Pipeline}
Die Kopplung klassischer \gls{gls-SIP}-Telefonie mit einem \gls{gls-LLM}-\gls{gls-API} bringt abweichende Audio-Formate zusammen: Im \gls{gls-PSTN} wird G.711-\gls{gls-PCM} mit 8 kHz und 8 Bit genutzt, wie im \gls{gls-RTP}-Profil beschrieben \parencite{rfc3551,ITU-T-G114}. Browserbasierte \gls{gls-WebRTC}-Pfade setzen hingegen auf den Opus-Codec mit interner 48 kHz Abtastung und variabler Bitrate \parencite{rfc6716}. Für die anschlie\ss ende Spracherkennung verlangen aktuelle \gls{gls-ASR}-Modelle wie Whisper \gls{gls-PCM16}-Audio bei 16 kHz \parencite{whisper_quant_2025}. Daraus entsteht eine zwingende Transcoding-Kette:

$\text{G.711 (8 kHz)} \xrightarrow{\text{\gls{gls-SIP} Bridge}} \text{Opus (48 kHz)} \xrightarrow{\text{Agent Worker}} \text{PCM16 (16 kHz)}$

Jede Stufe kann Verzögerung und mögliche Codec-Artefakte hinzufügen. Um den Overhead zu begrenzen, wird \gls{gls-VAD} eingesetzt: Sowohl das \gls{gls-RTP}-Profil (z.B. Annex B bei G.729) empfiehlt \gls{gls-VAD} zur Unterdrückung von Stillepaketen \parencite{rfc3551}, als auch der Opus-Pfad nutzt \gls{gls-VAD}-Metriken, um Sprachaktivität zu erkennen \parencite{rfc6716}. Durch das Ausblenden von Stille sinkt der Daten- und Token-Verbrauch spürbar, was sich insbesondere bei Cloud-\gls{gls-API} in geringeren Kosten niederschlägt.

Moderne \gls{gls-LLM}-\gls{gls-API} unterstützen Streaming-Antworten, sodass eine Ausgabe bereits gesendet wird, während das Modell noch restliche Tokens erzeugt. Whisper-Streaming liefert Live-Transkripte mit wenigen Sekunden Latenz und erlaubt Puffer-Feintuning für kontinuierliche Ausgaben \parencite{whisper_quant_2025}.

Damit sinkt die wahrgenommene Wartezeit, aber physische Verzögerungen bleiben: Die Kodierung selbst bringt Rahmen- und Vorlaufzeiten mit sich, und ein Jitter-Puffer in paketbasierten Netzen addiert im Mittel etwa die halbe erwartete Verzögerungsschwankung \parencite{rfc3550,ITU-T-G114}. Netzwerk-Laufzeiten und die Modellinferenz kommen zusätzlich hinzu, sodass trotz Streaming weiterhin spürbare Latenzquellen bestehen.

Es ist dabei wichtig zu unterscheiden: Während ITU-T G.114 für klassische Telefonie Einwegverzögerungen von rund 150 Millisekunden als unkritisch definiert \parencite{ITU-T-G114}, ist Whisper-Streaming für Echtzeitanwendungen konzipiert, obwohl es bei Live-Transkriptionen mit Verzögerungen von etwa 3,3 Sekunden arbeitet, wobei die Spracherkennung fortlaufend während des Sprechens erfolgt \parencite{whisper_quant_2025}. Daraus lässt sich ableiten, dass \gls{gls-KI}-basierte Dialoge Latenzen im unteren Sekundenbereich noch als natürlich erscheinen lassen, obwohl sie deutlich über den Limits der klassischen IP-Telefonie liegen.

Der zentrale Vorteil containerbasierter Architekturen wird deutlich, wenn \gls{gls-KI}-Dienste mit klassischen \gls{gls-SIP}-/\gls{gls-VoIP}-Systemen kombiniert werden: Während \gls{gls-SIP} weiterhin Signalisierung und Gesprächssteuerung übernimmt, können die Medienströme parallel zu \gls{gls-ASR}- und \gls{gls-LLM}-Services geleitet werden. Jede Komponente lässt sich dann unabhängig skalieren. Diese modulare Entkopplung ermöglicht es, Hochlast-Komponenten gezielt zu vergrö\ss ern, ohne das gesamte System zu überprovisionieren \parencite{Antonova2025,ReliabilityVirtualizedNetworks}.

\section{Rechtliche Aspekte und Datenschutz}

Da Sprachdaten personenbezogene Informationen enthalten, müssen vor allem in \gls{gls-KI}-basierten Telefonsystemen wichtige datenschutzrechtliche Grundlagen, insbesondere die \gls{gls-DSGVO}, zwingend berücksichtigt werden. Solche Systeme können nicht nur Identitätsinformationen, sondern auch sensible Zusatzdaten (z.B. Gesundheitsinformationen, Standorte, persönliche Präferenzen) nutzen und erfordern damit besonderen Schutz \parencite{Data2025_AIDataQuality}.

Kurzbefunde aus aktueller Forschung untermauern diese Sachlage und stellen heraus, dass die Verlagerung traditioneller \gls{gls-PBX}-Funktionen in containerbasierte Umgebungen zwar Flexibilität und Skalierbarkeit bringt, zugleich aber spezielle Sicherheits- und Orchestrierungsanforderungen (z.B. Netzwerksegmentierung, Schlüsselmanagement, kontrollierte Ingress-Punkte) notwendig macht, um personenbezogene Medienströme angemessen zu schützen \parencite{Antonova2025}. Darüber hinaus betonen Übersichtsarbeiten zur Zuverlässigkeit virtualisierter Netze, dass Virtualisierung die Interdependenzen zwischen Komponenten erhöht und damit neue Fehler und
Ausfallmodi einführt, die in Sicherheits- und
Datenschutzbewertungen explizit berücksichtigt werden müssen \parencite{ReliabilityVirtualizedNetworks}. Arbeiten zur Data-Governance weisen explizit auf die Bedeutung klarer Regelungen zur Datenlokalität, Retention-Politik und zur Nutzung von Subprozessoren hin, weil diese Punkte die rechtliche Bewertung von cloudbasierten \gls{gls-LLM}-Diensten ma\ss geblich beeinflussen \parencite{Data2025_AIDataQuality}.

Wesentliche rechtliche und ethische Prinzipien, die für\gls{gls-KI}-basierte Systeme gelten, basieren auf der \gls{gls-DSGVO} und
einschlägigen Data-Governance-Standards:

\noindent\textbf{Datenminimierung}
\begin{quote}
Es dürfen nur die unbedingt notwendigen Audiodaten und Metadaten verarbeitet werden.
\end{quote}
\noindent\textbf{Zweckbindung}
\begin{quote}
Aufzeichnung, Analyse und
Weiterverarbeitung sind auf den konkreten Zweck zu beschränken und vertraglich festzulegen.
\end{quote}

\noindent\textbf{Transparenz}
\begin{quote}
Betroffene sind über Art, Umfang und
Zweck der Verarbeitung zu informieren. Bei Telefonie ist dies in der Regel über Ansagen oder Datenschutzhinweise zu realisieren.
\end{quote}
\clearpage
\noindent\textbf{Integrität und Vertraulichkeit}
\begin{quote}
Technische Ma\ss nahmen (Verschlüsselung,
Zugriffskontrollen) müssen sicherstellen, dass Unbefugte keinen Zugriff erhalten.
\end{quote}

\noindent\textbf{Rechtmä\ss igkeit}
\begin{quote}
Verarbeitung muss auf einer Rechtsgrundlage beruhen (Einwilligung, Vertragserfüllung, berechtigtes Interesse, etc.).
\end{quote}

Diese Prinzipien lehnen sich an Artikel 5 der \gls{gls-DSGVO} an und werden durch aktuelle Data-Governance-Arbeiten konkretisiert \parencite{Data2025_AIDataQuality}.

\subsubsection{Konkrete Schutzma\ss nahmen}

Für den praktischen Schutz personenbezogener Mediendaten empfiehlt sich ein mehrschichtiges Vorgehen, das technische Härtung und organisatorische Prozesse kombiniert. Zunächst ist die strikte Trennung von Medien- und
Steuerkanälen zentral: Medienströme \gls{gls-RTP}/\gls{gls-SRTP} sollten in eigenen Netzwerkzonen oder in einem \gls{gls-VLAN} geführt werden. Dies sind logische Netzwerk-Segmente zur Isolation. In containerbasierten Umgebungen wie Kubernetes sind separate Namespaces mit expliziten \texttt{NetworkPolicy}-Regeln von Vorteil, um eine logische Trennung vorzunehmen. Ergänzend reduzieren dedizierte Ingress-Controller und restriktive Firewall-Policies die Angriffsfläche und begrenzen seitliche Übersprünge im Falle einer Kompromittierung. Dies entspricht den Empfehlungen für containerbasierte \gls{gls-PBX}-Bereitstellungen \parencite{Antonova2025}.

Ebenso wichtig ist ein zentrales und automatisiertes Geheimnis- und
Schlüsselmanagement: \gls{gls-API}-Keys, \gls{gls-TLS}-Zertifikate und Provider-Tokens dürfen nicht unverschlüsselt in Container-Images oder Quell-Repositories liegen. Stattdessen sollten Geheimnisse in einem gesicherten Speicher verwaltet, Kubernetes-Schlüssel nur verschlüsselt im Ruhezustand gehalten und Zugriffsrechte nach dem Least-Privilege-Prinzip vergeben werden. Kurzlebige Zugriffstokens und automatische Rotation vermindern das Risiko gestohlener Zugangsdaten. Antonova et al. heben genau diese Praxis für Asterisk-Deployments in Kubernetes hervor \parencite{Antonova2025}.

Observability und Logging müssen datenschutzgerecht gestaltet sein: Logs sind für Debugging, Monitoring und Forensik unverzichtbar, dürfen jedoch keine vollständigen Transkripte personenbezogener Gespräche in allgemeinen Log-Stores enthalten. Stattdessen empfiehlt sich die Ma\ss nahme des Maskierens oder Entfernens von personenbezogenen Informationen, ein separat verschlüsseltes Audit-Repository für sensible Ereignisse sowie klar definierte Aufbewahrungs- und Löschfristen, inklusive restriktiver Zugriffskontrollen. Diese technischen Praktiken sind notwendig, um Compliance-Anforderungen und Datenschutzverpflichtungen einhalten zu können \parencite{Data2025_AIDataQuality}.

Die Gewährleistung von Verfügbarkeit erfordert ebenfalls explizite Strategien: ReplicaSets und Auto-Scaling erhöhen die Ausfallsicherheit, lösen aber nicht automatisch Konsistenz- oder Persistenzfragen für zustandsbehaftete Daten (z.B. den \gls{gls-CDR} oder Ticket-Datenbanken). Deshalb sind regelmä\ss ige, getestete Backups (z.B. Snapshotting von Persistent-Volumes), dokumentierte Restore-Prozeduren und konsistente Sicherungsstrategien unabdingbar. Die Literatur zur Zuverlässigkeit virtualisierter Netze beschreibt gängige Modellierungsansätze und betont Testbarkeit und Wiederherstellbarkeit als zentrale Anforderungen \parencite{ReliabilityVirtualizedNetworks}.

Schlie\ss lich müssen technische Ma\ss nahmen durch betriebliche Verfahren ergänzt werden: Privacy-Impact-Assessments, regelmä\ss ige Sicherheits- und Compliance-Audits, dokumentierte Incident-Response-Pläne sowie zielgerichtete Schulungen für Entwickler und Operatoren schaffen die organisatorische Grundlage, damit die technischen Kontrollen wirksam bleiben. Aktuelle Data-Governance-Publikationen und Arbeiten zu Governance-Fragen bei \gls{gls-LLM}-Pipelines empfehlen genau diese Kombination aus Technik und Prozesssteuerung \parencite{Data2025_AIDataQuality,ssrn5348245}.

Neben Privacy-Impact-Assessments empfiehlt die Literatur die Durchführung von \gls{gls-EIA}, um ethische Risiken und Auswirkungen systematisch zu identifizieren und zu steuern. Zentrale Empfehlungen umfassen zudem die Benennung spezifischer Verantwortlicher, die kontinuierliche Überwachung von Schutzmaßnahmen, die Verbreitung von Best Practices sowie gezielte Schulungen für Mitarbeitende. Insbesondere für Hochrisiko-\gls{gls-KI}-Systeme werden Transparenz und Erklärbarkeit als zentrale Anforderungen hervorgehoben, um Vertraulichkeit und Nachvollziehbarkeit zu gewährleisten \parencite{Schmiedchen2026KI}.

\subsubsection{Datenschutzrisiken bei Einsatz von \gls{gls-LLM}}

Viele \gls{gls-LLM}- und Sprach-\gls{gls-API} werden als Cloud-Services angeboten. Dies vereinfacht eine Integration und den Betrieb, bringt jedoch spezifische Risiken mit sich, die einer systematischen Analyse bedürfen. Nach dem JusGov-Rahmenwerk gelten Aspekte wie Notwendigkeit, Proportionalität und Risikobeurteilung als zentral für eine \gls{gls-DPIA} nach Artikel 35 der \gls{gls-GDPR} \parencite{ssrn5348245}.\\Im Kontext von Cloud-\gls{gls-LLM}-Diensten sind folgende Kernrisiken relevant:
\clearpage
\begin{description}
\item[Datenübermittlung in Drittstaaten]
Es besteht das Risiko, dass Audiodaten und daraus erzeugte Transkripte bei der Nutzung cloudbasierter \gls{gls-KI}-Dienste au\ss erhalb der EU verarbeitet oder durch Subprozessoren in Drittstaaten weitergegeben werden. Fehlen geeignete Garantien, wie etwa Zertifizierungen nach \gls{gls-ISO}~27001 oder vertragliche Absicherungen, ergeben sich erhöhte datenschutzrechtliche Risiken im Sinne der \gls{gls-DSGVO} \parencite{Data2025_AIDataQuality}.

\item[Begrenzte Einsicht und Auditierbarkeit]
Die interne Verarbeitungspipeline eines Anbieters entzieht sich direkter Kontrolle. Risiken entstehen durch mangelnde Transparenz in der Speicherung, Löschung und Nachvollziehbarkeit von personenbezogenen Daten \parencite{ssrn5348245}.

\item[Unkontrollierte Datenspeicherung]
Anbieter können Nutzereingaben (Audio, Transkripte, Metadaten) zu Monitoring- oder Qualitätszwecken protokollieren oder über vereinbarte Zeiträume hinaus speichern. Ohne explizite vertragliche Aufbewahrungsfristen besteht das Risiko unerwünschter Persistenz und potenzieller Zweckentfremdung \parencite{Data2025_AIDataQuality}.

\item[Bias und Diskriminierende Muster]
\gls{gls-LLM} können durch Training auf historisch verzerrten Daten diskriminierende Outputs erzeugen und damit gegen das Trustworthy-\gls{gls-AI}-Prinzip der Fairness versto\ss en \parencite{ssrn5348245}.
\end{description}

In der nachfolgenden Aufstellung sind wesentliche Risiken und deren empfohlene Schutzma\ss nahmen tabellarisch zusammengefasst:

\begin{table}[H]
  \centering
  \caption[Risiken von Cloud-\gls{gls-LLM}-Diensten]{Risiken von Cloud-\gls{gls-LLM}-Diensten und empfohlene Schutzma\ss nahmen (Angelehnt an: \textcite{Data2025_AIDataQuality}, \textcite{ssrn5348245})}
  \label{tab:llm_risks}
  {\small
  \setlength{\extrarowheight}{4pt}
  \begin{tabular}{@{} >{\raggedright\arraybackslash}p{0.48\textwidth} >{\raggedright\arraybackslash}p{0.48\textwidth} @{} }
    \toprule
    \textbf{Risiko} & \textbf{Ma\ss nahmen} \\
    \midrule
    Datenübermittlung in Drittstaaten & \gls{gls-SCC}, Datenminimierung, Verschlüsselung vor Übertragung, EU-Datenorte vertraglich festlegen \\
    \addlinespace[2pt]
    Begrenzte Einsicht/Auditierbarkeit & On-premise/Edge-Optionen für kritische Workloads, Lösch-/Audit-Rechte im \gls{gls-DPA}, transparente Subprozessoren-Listen \\
    \addlinespace[2pt]
    Unkontrollierte Datenspeicherung & Explizite Aufbewahrungsfristen im \gls{gls-DPA}, Log-Redaction (Schwärzung sensibler Daten), Metadaten-Minimierung, regelmä\ss ige Kontrollaudits \\
    \addlinespace[2pt]
    Bias und Diskriminierung & Datenqualitäts-Prüfungen, multidisziplinäre Audits (IT/Recht/Ethik), Überwachung auf diskriminierende Outputs, \gls{gls-DPIA}-Dokumentation \\
  \end{tabular}
  }
  \\[1em]
  \footnotesize{\textit{Quellen:} \textcite{ssrn5348245}, \textcite{Data2025_AIDataQuality}}
\end{table}

Neben der Vereinbarung über ein \gls{gls-DPA} sollte vor einem Produktivbetrieb wesentliche Ma\ss nahmen, beispielsweise in Form einer formalen \gls{gls-DPIA}, durchgeführt werden. Die zentralen Aspekte einer formalen \gls{gls-DPIA} sind:
\begin{itemize}
  \item \textbf{Notwendigkeit prüfen}: Ist Cloud-\gls{gls-LLM} wirklich erforderlich oder gibt es weniger invasive Alternativen?
  \item \textbf{Stakeholder-Konsultation}: Nach Möglichkeit sollten Betroffene (Anrufer, Mitarbeiter) in die Bewertung einbezogen werden. Dies orientiert sich an Artikel 35 Absatz 9 \gls{gls-GDPR}.
  \item \textbf{Fortbestehende Risiken}: Falls nach einer Mitigation hohe Risiken bestehen, ist die zuständige Datenschutzbehörde zu konsultieren. Bezugnehmend auf Artikel 36 \gls{gls-GDPR}.
\end{itemize}

Eine \gls{gls-DPIA} ist als dynamisches Dokument zu verstehen und sollte bei wesentlichen Änderungen an der Verarbeitung oder dem technischen Setup aktualisiert werden, um eine kontinuierliche Compliance zu gewährleisten \parencite{ssrn5348245}.

Die dargestellten Grundlagen zeigen, dass containerbasierte Deployment-Modelle, standardisierte \gls{gls-VoIP}-Protokolle und datenschutzkonforme Verarbeitungsprozesse die technische und rechtliche Basis für ein \gls{gls-KI}-basiertes Telefonsystem bilden.
